% ============================================================
% showcase-paper.tex — harryopo-paper.cls 全功能展示模板
% 覆盖：文字样张、层级架构、中英文混排、公式、表格、
%        图像、代码/算法、定理、引用、列表等全部特性
% 编译：xelatex(3次)
% ============================================================
\documentclass{harryopo-paper}

% --- 论文信息 ---
\title{harryopo LaTeX 论文模板全功能展示\\——从基础排版到高级特性}
\author{模板作者\quad 测试用户}
\date{2026年6月21日}
\fundproject{基金项目：模板开发专项（No. harryopo20260001）}

\begin{document}

\maketitle

\begin{abstract}
本文档是学术论文模板（\texttt{harryopo-paper.cls}）的全功能展示文件，旨在帮助作者快速上手使用本模板。文档系统性地展示了模板的全部排版能力，涵盖：中文与英文混排文本、多级标题层级架构（节/小节/括号细目）、行内公式与行间公式（\texttt{unicode-math} + XITS Math）、表格（基础三线表/列宽自适应/多行合并）、图片插入与子图排版、算法伪代码（\texttt{algorithm2e}）、Python/Bash 代码高亮（\texttt{listings}）、定理与定义数学环境、文献引用与上标引用、列表环境（无序/有序/嵌套）、以及基金信息与作者简介等辅助元素。作者可直接复制本文档中的代码片段到自己的论文中使用，实现零门槛上手。
\end{abstract}

\keywords{LaTeX模板；学术论文；全功能展示；排版规范；harryopo}

% ============================================================
\section{文字与段落排版}
\label{sec:text}

\subsection{中文正文混排英文}

中文段落中自然嵌入英文词汇（如 \emph{LaTeX}、\emph{Convolutional Neural Network}、\emph{Backpropagation}）以及阿拉伯数字（1，2，3，…）时，模板自动使用 \emph{XITS} 字体族排版英文与数字。XITS 是与 XITS Math 数学字体配套的 Times 风格衬线字体，确保中西文混排时视觉风格统一、基线对齐自然。

以下是一段较长的混排示例文本，展示了模板在中英文混排场景下的实际效果：

The Transformer architecture, introduced by Vaswani et al.\ in the seminal paper \textquotedblleft Attention Is All You Need\textquotedblright{} (NeurIPS 2017), has fundamentally reshaped the landscape of deep learning research. 凭借其自注意力机制（Self-Attention）与多头设计（Multi-Head Attention），Transformer 不仅在自然语言处理领域取得了 GPT-4、LLaMA 等里程碑式成果，还被成功迁移到计算机视觉（Vision Transformer, ViT）、语音识别（Conformer）和生物计算（AlphaFold）等广泛领域。2025 年以来，混合专家模型（Mixture of Experts, MoE）与状态空间模型（State Space Model, SSM）进一步推动了架构演进，参数规模突破万亿量级，标志着通用人工智能时代的加速到来。

{技术演进路线如下：CNN（2012）→ RNN/LSTM（2014）→ ResNet（2016）→ Transformer（2017）→ BERT/GPT（2018）→ ViT（2020）→ GPT-4（2023）→ 多模态大模型（2024--2025）}

\subsection{多级标题层级架构}

本模板提供了三个标准层级的编号标题（\verb|\section| / \verb|\subsection| / \verb|\subsubsection|）以及无编号括号式细目标题（\verb|\subhead|），完整覆盖学术论文的所有层级需求。

\subsubsection{三级编号标题（本行为 subsubsection 示例）}

三级编号标题使用方正楷体，左缩进 2em，字体大小与正文一致但颜色略深。适用于需要在二级标题下进一步细分、但又不想增加编号深度的场景。

\subhead{括号式细目标题}

括号式细目标题是一个无编号的黑体标注，以中文括号包裹标题文字，后接正文。适用于更细粒度的逻辑分割，避免产生四层以上的编号嵌套。

\subhead{另一个括号细目}两个括号细目之间通过段间距自然分隔，无需手动添加空行即可区分逻辑边界。这种设计既保持了版面紧凑度（不浪费垂直空间），又提供了清晰的视觉提示。

\subsection{文字效果与强调}

模板提供以下常用文字效果命令：\textbf{粗体强调}、\textit{斜体（英文语境）}、\texttt{等宽字体}、\emph{强调文本}、\underline{下划线文本}。中文楷体可通过 \verb|\fzkt| 命令调用：{\fzkt 这一句使用方正楷体显示}；中文仿宋为 {\fzfs 方正仿宋}；中文黑体为 {\fzht 方正黑体}。

\subsection{列表环境}

模板已配置紧凑型列表（\verb|enumitem|），支持无序列表、有序列表和嵌套列表：

\begin{itemize}
    \item 第一层无序列表项：模板使用紧凑间距
    \begin{itemize}
        \item 第二层无序子项：自动缩进
        \item 另一子项：间距控制良好
    \end{itemize}
    \item 回到第一层
\end{itemize}

\begin{enumerate}
    \item 第一项：有序列表
    \item 第二项：带子列表
    \begin{enumerate}
        \item 子项 (a)
        \item 子项 (b)
    \end{enumerate}
    \item 第三项
\end{enumerate}

% ============================================================
\section{数学公式排版}
\label{sec:math}

本模板使用 \texttt{unicode-math} 包配合 \emph{XITS Math} 字体引擎，支持完整的 Unicode 数学符号体系。以下按公式类型逐一展示。

\subsection{行内公式}

行内公式与正文在同一行内排列，如质能方程 $E = mc^{2}$、高斯积分 $\int_{-\infty}^{\infty} e^{-x^{2}}\,dx = \sqrt{\pi}$、集合定义 $\forall\varepsilon>0,\ \exists\delta>0$ 等。带上下标的符号如 $\mathbf{W}_{ij}^{(l)}$ 或希腊字母如 $\alpha, \beta, \gamma, \theta, \lambda, \mu, \sigma, \phi, \omega$ 均可自然排版。

行内分数使用 \verb|\nicefrac| 或斜线形式：$x = 1/2$，复杂情况建议使用行间公式。

\subsection{单行行间公式}

\begin{equation}
f(x) = \frac{1}{\sigma\sqrt{2\pi}} \exp\left(-\frac{(x - \mu)^{2}}{2\sigma^{2}}\right)
\label{eq:gaussian}
\end{equation}

式\eqref{eq:gaussian}是正态分布的概率密度函数（equation 环境）。注意本模板将公式编号统一为 `(节号.公式序号)' 格式。

\subsection{多行对齐公式}

\begin{align}
\tilde{D}^{-\frac{1}{2}} \tilde{A} \tilde{D}^{-\frac{1}{2}} \quad\text{其中}\quad
\tilde{A} &= A + I_{N} \label{eq:gcn_pre}\\
H^{(l+1)} &= \sigma\!\left( \tilde{D}^{-\frac{1}{2}} \tilde{A} \tilde{D}^{-\frac{1}{2}} H^{(l)} W^{(l)} \right) \label{eq:gcn}
\end{align}

式\eqref{eq:gcn_pre}--\eqref{eq:gcn}为图卷积层的前向传播定义（align 环境），每行可独立编号并通过 \verb|\label| 交叉引用。

\subsection{分段函数与矩阵}

\begin{equation}
\operatorname{ReLU}(x) = \begin{cases}
x, & \text{若 } x > 0 \\
0, & \text{若 } x \leq 0
\end{cases}
\label{eq:relu}
\end{equation}

\begin{equation}
\mathbf{A} = \begin{pmatrix}
a_{11} & a_{12} & \cdots & a_{1n} \\
a_{21} & a_{22} & \cdots & a_{2n} \\
\vdots & \vdots & \ddots & \vdots \\
a_{m1} & a_{m2} & \cdots & a_{mn}
\end{pmatrix}
\label{eq:matrix}
\end{equation}

\subsection{多行条件推导}

\begin{align}
\mathcal{L}_{\text{total}} &= \mathcal{L}_{\text{cls}} + \lambda_{1} \mathcal{L}_{\text{reg}} + \lambda_{2} \mathcal{L}_{\text{aux}} \label{eq:loss}\\
\mathcal{L}_{\text{cls}} &= -\frac{1}{N}\sum_{i=1}^{N} \sum_{c=1}^{C} y_{ic} \log \hat{y}_{ic} \label{eq:ce}\\
\mathcal{L}_{\text{reg}} &= \sum_{l=1}^{L} \|\mathbf{W}^{(l)}\|_{F}^{2} \label{eq:l2}
\end{align}

其中 $\lambda_{1}, \lambda_{2} \in \mathbb{R}^{+}$ 为平衡超参数，$\|\cdot\|_{F}$ 为 Frobenius 范数。

\subsection{定理与证明环境}

模板预置了 \verb|definition|、\verb|theorem|、\verb|lemma|、\verb|corollary|、\verb|example| 等定理类环境（均为 \texttt{amsthm} 定义）。

\begin{definition}[卷积运算]
给定输入信号 $x(t)$ 与卷积核 $k(t)$，其离散卷积定义为：
\begin{equation}
y[n] = (x * k)[n] = \sum_{m=-\infty}^{\infty} x[m] \cdot k[n - m]
\end{equation}
\end{definition}

\begin{theorem}[通用逼近定理]
对于任意连续函数 $f$ 定义在紧致子集 $K \subset \mathbb{R}^{n}$ 上，以及任意 $\varepsilon > 0$，存在一个具有单隐藏层的前馈神经网络 $g$，使得：
\begin{equation}
\sup_{\mathbf{x} \in K} \|f(\mathbf{x}) - g(\mathbf{x})\| < \varepsilon
\end{equation}
\end{theorem}

\begin{lemma}
设 $\mathcal{N}(\mu, \sigma^{2})$ 为正态分布，则其矩母函数为
$M(t) = \exp(\mu t + \frac{1}{2}\sigma^{2}t^{2})$。
\end{lemma}

\begin{proof}
由矩母函数定义，$M(t) = \mathbb{E}[e^{tX}] = \int_{-\infty}^{\infty} e^{tx} \frac{1}{\sigma\sqrt{2\pi}} e^{-(x-\mu)^{2}/(2\sigma^{2})}\,dx$。配方后积分可得结论。
\end{proof}

\begin{corollary}
正态分布的任意阶矩可由 $\mu$ 和 $\sigma^{2}$ 表示。
\end{corollary}

% ============================================================
\section{表格排版}
\label{sec:table}

\subsection{基础三线表}

表\ref{tab:basic}展示基础的三线表（\texttt{booktabs}）风格，这是学术论文中最常用的表格形式。

\begin{table}[htbp]
\centering
\begin{tabular}{lccc}
\toprule
\textbf{框架} & \textbf{发布年份} & \textbf{动态图} & \textbf{生态成熟度} \\
\midrule
TensorFlow   & 2015 & \checkmark & 很高 \\
PyTorch      & 2016 & \checkmark & 很高 \\
JAX          & 2018 & \checkmark & 中等 \\
PaddlePaddle & 2016 & \checkmark & 较高 \\
MindSpore    & 2020 & \checkmark & 中等 \\
\bottomrule
\end{tabular}
\caption{主流深度学习框架特性对比}
\label{tab:basic}
\end{table}

\subsection{列宽自适应表格}

表\ref{tab:tabularx}使用 \texttt{tabularx} 环境实现列宽自适应，适用于单元格内容较长的场景。

\begin{table}[htbp]
\centering
\begin{tabularx}{0.92\textwidth}{lXX}
\toprule
\textbf{优化器} & \textbf{核心思想} & \textbf{适用场景} \\
\midrule
SGD + Momentum & 动量累积历史梯度方向，加速收敛并抑制震荡 & 图像分类、目标检测等 CV 任务 \\
Adam & 结合动量与自适应学习率，对稀疏梯度友好 & NLP、Transformer 训练 \\
AdamW & 解耦权重衰减与自适应学习率，正则化更强 & 大模型预训练、微调 \\
LAMB & 分层自适应学习率，支持超大 batch size & 分布式大规模训练 \\
\bottomrule
\end{tabularx}
\caption{优化器配置与适用场景（\texttt{tabularx} 自适应列宽）}
\label{tab:tabularx}
\end{table}

\subsection{多行合并表格}

表\ref{tab:multirow}展示 \texttt{multirow} 和 \texttt{multicolumn} 的组合使用。

\begin{table}[htbp]
\centering
\begin{tabular}{lccc}
\toprule
\multirow{2}{*}{\textbf{数据集}} & \multirow{2}{*}{\textbf{样本数}} & \multicolumn{2}{c}{\textbf{图像尺寸}} \\
\cmidrule(lr){3-4}
   &  & \textbf{训练} & \textbf{测试} \\
\midrule
CIFAR-10    & 60,000  & $32\times32$ & $32\times32$ \\
ImageNet-1K & 1,281,167 & $224\times224$ & $224\times224$ \\
COCO        & 118,287 & 可变 & 可变 \\
\bottomrule
\end{tabular}
\caption{实验数据集统计信息（含多行合并）}
\label{tab:multirow}
\end{table}

% ============================================================
\section{图片插入}
\label{sec:image}

\subsection{单图插入}

图\ref{fig:demo}使用 \texttt{figure} 环境居中排版，配合 \texttt{caption} 包生成图题。本示例用 TikZ 绘制示意图形以替代外部图片文件。

\begin{figure}[htbp]
\centering
\begin{tikzpicture}[scale=1.2]
    \fill[MainColor!15]  (0,0)  rectangle (2,1.5);
    \fill[SubColor!25]   (1,0)  rectangle (3,1.5);
    % 网络层级标注
    \draw[->,thick,MainColor] (-1.5,0.75) -- (-0.5,0.75) node[midway,above,font=\footnotesize] {输入};
    \draw[->,thick,SubColor]  (3.5,0.75)  -- (4.5,0.75)  node[midway,above,font=\footnotesize] {输出};
    \node[font=\small\bfseries,fill=white,inner sep=2pt] at (1,0.75) {Encoder};
    \node[font=\small\bfseries,fill=white,inner sep=2pt] at (2,0.75) {Decoder};
    % 下方说明
    \node[font=\footnotesize\fzkt,MainColor!70!black] at (1.5,-0.6) {Transformer 编码器-解码器结构示意};
\end{tikzpicture}
\caption{Transformer 编码器-解码器架构示意（TikZ 内联绘制）}
\label{fig:demo}
\end{figure}

\subsection{子图排版}

图\ref{fig:subfig}展示使用 \texttt{subcaption} 包的并列子图排版。

\begin{figure}[htbp]
\centering
\begin{subfigure}{0.45\textwidth}
\centering
\begin{tikzpicture}[scale=0.9]
    \fill[blue!10] (0,0) rectangle (2,2);
    \draw[->,blue!60,thick] (0.5,1) -- (1.5,1.5);
    \node[font=\footnotesize] at (1,0.3) {(a) 线性可分};
\end{tikzpicture}
\caption{线性可分}
\end{subfigure}
\hfill
\begin{subfigure}{0.45\textwidth}
\centering
\begin{tikzpicture}[scale=0.9]
    \fill[red!10] (0,0) rectangle (2,2);
    \draw[red!60,thick] plot[smooth] coordinates {(0.3,0.5)(0.8,1)(1.2,1.2)(1.7,1.3)};
    \node[font=\footnotesize] at (1,0.3) {(b) 非线性可分};
\end{tikzpicture}
\caption{非线性可分}
\end{subfigure}
\caption{分类任务的两种典型场景（子图并列）}
\label{fig:subfig}
\end{figure}

% ============================================================
\section{代码与算法}
\label{sec:code}

\subsection{Python 代码高亮}

模板预设了 GitHub 风格的 Python 代码高亮方案（\texttt{pystyle}），支持行号和语法着色。

\begin{lstlisting}[style=pystyle, caption={Python 示例：SGD 优化器实现}, label=lst:python]
import torch
from torch import nn, optim

class SimpleCNN(nn.Module):
    def __init__(self, num_classes=10):
        super().__init__()
        self.conv = nn.Sequential(
            nn.Conv2d(3, 32, 3, padding=1), nn.ReLU(),
            nn.MaxPool2d(2),
            nn.Conv2d(32, 64, 3, padding=1), nn.ReLU(),
            nn.MaxPool2d(2),
        )
        self.fc = nn.Linear(64 * 8 * 8, num_classes)

    def forward(self, x):
        x = self.conv(x)
        return self.fc(x.view(x.size(0), -1))
\end{lstlisting}

\subsection{Bash 终端命令}

\begin{lstlisting}[style=bashstyle, caption={训练命令示例}, label=lst:bash]
# 安装依赖
pip install torch torchvision numpy matplotlib tqdm

# 单卡训练
python train.py --epochs 100 --batch_size 128 --lr 0.01 \
    --dataset cifar10 --arch resnet18

# 多卡分布式训练
torchrun --nproc_per_node=4 train.py --distributed \
    --dataset imagenet --arch vit_base
\end{lstlisting}

\subsection{算法伪代码}

本模板使用 \texttt{algorithm2e} 宏包排版算法伪代码，中文关键字已适配（算法/输入/输出/函数）。

\begin{algorithm}[htbp]
\caption{梯度下降优化算法}
\label{alg:sgd}
\KwIn{训练集 $\mathcal{D}$，学习率 $\eta$，最大迭代次数 $T$}
\KwOut{模型参数 $\theta$}
\SetKwFunction{GradLoss}{Grad}
\BlankLine
初始化参数 $\theta_0$\;
\For{$t = 1$ \KwTo $T$}{
    从 $\mathcal{D}$ 中采样 mini-batch $\mathcal{B}$\;
    计算梯度 $\mathbf{g}_t \leftarrow \frac{1}{|\mathcal{B}|}\sum_{(\mathbf{x},y) \in \mathcal{B}} \GradLoss{\mathbf{x}, y, \theta_{t-1}}$\;
    更新参数 $\theta_t \leftarrow \theta_{t-1} - \eta \cdot \mathbf{g}_t$\;
    \If{$\|\mathbf{g}_t\| < \varepsilon$}{
        收敛，退出循环\;
    }
}
\Return{$\theta_T$}
\end{algorithm}

% ============================================================
\section{参考文献引用}
\label{sec:cite}

\subsection{引用方式}

本模板支持常规引用 \verb|\cite{key}| 和上标引用 \verb|\upcite{key}| 两种方式。以下为引用示例：

\begin{itemize}
    \item 深度学习综述：\cite{lecun2015deep}
    \item ResNet 残差网络：\cite{he2016deep}
    \item Attention 机制原始论文：\cite{vaswani2017attention}
    \item ImageNet 数据集：\cite{deng2009imagenet}
    \item Batch Normalization：\cite{ioffe2015batch}
    \item Dropout 正则化：\cite{srivastava2014dropout}
    \item 中文引用：\cite{zhou2016machine}
\end{itemize}

上标引用示例：深度学习技术\upcite{lecun2015deep}已广泛应用于计算机视觉领域\upcite{he2016deep}，Transformer 架构\upcite{vaswani2017attention}则推动了自然语言处理的发展。

\subsection{交叉引用}

跨章节/图/表/公式的交叉引用通过 \verb|\ref| / \verb|\eqref| 命令实现。例如：文字排版详见第\ref{sec:text}节；公式排版详见第\ref{sec:math}节；单行公式\eqref{eq:gaussian}为高斯分布；图\ref{fig:demo}为 Transformer 架构示意；表\ref{tab:basic}为框架特性对比；算法\ref{alg:sgd}为梯度下降伪代码。

% ============================================================
\section{模板选项与自定义}
\label{sec:options}

\subsection{文档类选项}

\texttt{harryopo-paper.cls} 支持以下可选参数：

\begin{itemize}
    \item \verb|twocolumn| — 启用双栏排版模式
    \item \verb|dark| — 切换为深紫色主题配色
\end{itemize}

例如 \verb|\documentclass[twocolumn]{harryopo-paper}| 将启用双栏论文模板，\verb|\documentclass[dark]{harryopo-paper}| 启用深紫色主题。双栏模式下的完整示例请参阅 \texttt{example-paper-twocolumn.tex}。

\subsection{常用命令速查}

表\ref{tab:commands}汇总了本模板的核心使用命令。

\begin{table}[htbp]
\centering
\begin{tabular}{lp{7cm}}
\toprule
\textbf{命令} & \textbf{说明} \\
\midrule
\verb|\fzkt| / \verb|\fzfs| / \verb|\fzht| & 方正楷体 / 仿宋 / 黑体 \\
\verb|\fzxb| / \verb|\fzdbs| & 方正小标宋 / 大标宋 \\
\verb|\subhead{标题}| & 括号式细目小标题 \\
\verb|\maketitle| & 生成论文标题区 \\
\verb|\keywords{...}| & 关键词（分号分隔） \\
\verb|\upcite{key}| & 上标引用 \\
\verb|\printfunding| & 打印基金信息 \\
\verb|\printauthorbio| & 打印作者简介 \\
\bottomrule
\end{tabular}
\caption{模板核心命令速查}
\label{tab:commands}
\end{table}

% ==== 基金信息 ====
\printfunding

% --- 参考文献 ---
\begin{thebibliography}{9}

\bibitem{lecun2015deep} LeCun Y, Bengio Y, Hinton G. Deep learning[J]. \emph{Nature}, 2015, 521(7553): 436--444.

\bibitem{he2016deep} He K, Zhang X, Ren S, et al. Deep residual learning for image recognition[C]. \emph{CVPR}, 2016: 770--778.

\bibitem{vaswani2017attention} Vaswani A, Shazeer N, Parmar N, et al. Attention is all you need[C]. \emph{NeurIPS}, 2017: 5998--6008.

\bibitem{deng2009imagenet} Deng J, Dong W, Socher R, et al. ImageNet: A large-scale hierarchical image database[C]. \emph{CVPR}, 2009: 248--255.

\bibitem{ioffe2015batch} Ioffe S, Szegedy C. Batch normalization: Accelerating deep network training by reducing internal covariate shift[C]. \emph{ICML}, 2015: 448--456.

\bibitem{srivastava2014dropout} Srivastava N, Hinton G, Krizhevsky A, et al. Dropout: A simple way to prevent neural networks from overfitting[J]. \emph{JMLR}, 2014, 15(1): 1929--1958.

\bibitem{zhou2016machine} 周志华. 机器学习[M]. 北京: 清华大学出版社, 2016.

\end{thebibliography}

\end{document}
